\section{Regional embeddings in uniform physical dimension}
\label{sec:peps_region_operator_embedding}

Let $V$ be finite, let $q\in\N$, and set $A_q=\{0,\ldots,q-1\}$,
with $A_0=\varnothing$. For $S\subseteq V$, write $X_S=A_q^S$ and
$\mathcal H_S=\C^{X_S}$. Fix a region $R\subseteq V$, its inclusion
$j_R:R\hookrightarrow V$, and an arbitrary matrix
$H\in\operatorname{End}(\mathcal H_R)$.
The restriction map identifies $X_V$ with $X_R\times X_{V\setminus R}$;
let $L_R(H)$ denote $H\otimes I_{V\setminus R}$ in the $X_V$ coordinates.
For a site inclusion $j$, let $E_j(H)$ act by $H$ on its image and by the
identity on the remaining sites.
% Original comparison of existing regional and circuit definitions, supporting
% the operator conventions of 05-frames.tex, sec:frames, eq:hole-projector,
% eq:hole-encoder, lines 13--97, at adc7f1241b42e322a6451854ab7e4b4c146bf78a.
% This is not a source-labelled frame, reset-ownership, or approximation theorem.

\begin{theorem}[Equality of regional and site embeddings]
    \label{thm:peps_region_site_embedding}
    \lean{TNLean.PEPS.agreeOff_region_iff,
        TNLean.PEPS.regionLocalTerm_eq_embedOp}
    \leanok
    For every $\sigma,\tau\in X_V$,
    \begin{align}
        L_R(H)_{\sigma,\tau}\label{eq:peps_region_embedding_entries}
         & =(E_{j_R}(H))_{\sigma,\tau}
        =\begin{cases}
             H_{\sigma|_R,\tau|_R},
                & \sigma|_{V\setminus R}=\tau|_{V\setminus R}, \\
             0, & \text{otherwise}.
         \end{cases}
    \end{align}
\end{theorem}
\begin{proof}
    \leanok
    Since $j_R(R)=R$, agreement outside the image of $j_R$ is precisely
    equality of the restrictions to $V\setminus R$.
    The complementary identity matrix contributes the Kronecker delta
    $\delta_{\sigma|_{V\setminus R},\tau|_{V\setminus R}}$.
    Multiplying it by $H_{\sigma|_R,\tau|_R}$ gives
    (\ref{eq:peps_region_embedding_entries}) for both operators.
\end{proof}

\begin{theorem}[Full-region coordinate identification]
    \label{thm:peps_region_embedding_coordinates}
    \lean{TNLean.PEPS.dependentRegionOperatorLift_eq_reindex_embedOp}
    \leanok
    Let $\widetilde X_V=\prod_{v\in V}A_q$ and
    $\theta_q:X_V\to\widetilde X_V$ be the bijection
    $\theta_q(\sigma)=(\sigma(v))_{v\in V}$.
    Write $D_R(H)$ for $H\otimes I_{V\setminus R}$ after splitting this
    indexed product into its regional and complementary factors.
    For every $\alpha,\beta\in\widetilde X_V$,
    \begin{align}
        D_R(H)_{\alpha,\beta}
         & =(E_{j_R}(H))_{\theta_q^{-1}(\alpha),\theta_q^{-1}(\beta)}.
        \label{eq:peps_region_embedding_coordinates}
    \end{align}
\end{theorem}
\begin{proof}
    \leanok
    The restrictions of $\theta_q^{-1}(\alpha)$ and
    $\theta_q^{-1}(\beta)$ select the same coordinates as the two indexed
    product factors. Both sides are $H$ at the regional coordinates when
    the complementary coordinates agree, and zero otherwise.
\end{proof}

\begin{remark}[Empty sets and zero physical dimension]
    If $R=\varnothing$, then $X_R$ has one element and $L_R([c])=cI$.
    If $R=V$, there is no complementary agreement condition.
    If $q=1$, all configuration sets have one element, so every regional
    operator extends as a scalar identity.
    If $q=0$ and $V\ne\varnothing$, then $X_V=\varnothing$ and the
    resulting operator is zero. If $V=\varnothing$, then $X_V$ still has
    one element, including when $q=0$, and its scalar entry is retained.
\end{remark}

\begin{example}[Complex entries and the complementary identity]
    Take $V=\{0,1\}$, $R=\{0\}$, $q=2$, and
    $H=\bigl(\begin{smallmatrix}0&i\\2+i&0\end{smallmatrix}\bigr)$.
    Formula~(\ref{eq:peps_region_embedding_entries}) gives
    $L_R(H)_{(0,0),(1,0)}=i$,
    $L_R(H)_{(1,0),(0,0)}=2+i$, and
    $L_R(H)_{(0,0),(1,1)}=0$.
\end{example}
